\setchapterpreamble[u]{\margintoc}
\chapter{Jacobians, Hessians, and Taylor Approximation}
\labch{jacobian-hessian}

\sources{Hass, Heil and Weir, \emph{Thomas' Calculus}, 14th ed.\ \cite{thomas2018calculus},
Chapters~14--15; MIT 18.02 \cite{ocw1801}.}

The gradient handles scalar-valued functions. Vector-valued functions need a
matrix of derivatives, and second-order behaviour needs another. Both are
already familiar objects from Part~I wearing new labels.

\section{The Jacobian}
\labsec{jacobian}

\begin{definition}
For \(\vect{f}:\R^{n}\to\R^{m}\), the \emph{Jacobian} at \(\vect{x}\) is
\(J\in\R^{m\times n}\) with entries
\[
	J_{ij}=\frac{\partial f_i}{\partial x_j}.
\]
\end{definition}

Row \(i\) of \(J\) is the gradient of the \(i\)th output. The first-order
approximation is
\[
	\vect{f}(\vect{x}+\vect{h})=\vect{f}(\vect{x})+J\vect{h}+o(\lVert\vect{h}\rVert),
\]
so the Jacobian is the matrix of the linear map that best imitates \(\vect{f}\)
near \(\vect{x}\). When \(m=1\) it is \(\nabla f\tp\), a single row.

\marginnote{Note the transposition. The gradient is a column by convention, the
Jacobian of a scalar function a row. Keeping this straight is the source of most
sign-free errors in gradient derivations.}

\begin{example}
An affine layer \(\vect{f}(\vect{x})=W\vect{x}+\vect{b}\) has \(J=W\)
everywhere: it is already linear, so its best linear approximation is itself. An
elementwise activation \(\vect{f}(\vect{x})=\bigl(\phi(x_1),\ldots,\phi(x_n)\bigr)\)
has \(J=\operatorname{diag}\bigl(\phi'(x_1),\ldots,\phi'(x_n)\bigr)\).
\end{example}

\section{The chain rule in matrix form}
\labsec{matrix-chain-rule}

\begin{theorem}
If \(\vect{h}=\vect{g}\circ\vect{f}\) with \(\vect{f}:\R^{n}\to\R^{p}\) and
\(\vect{g}:\R^{p}\to\R^{m}\), then
\[
	J_{\vect{h}}(\vect{x})=J_{\vect{g}}\bigl(\vect{f}(\vect{x})\bigr)\,J_{\vect{f}}(\vect{x}).
\]
\end{theorem}

The scalar chain rule of \refsec{chain-rule} is the case \(m=p=n=1\). The
matrix version says a composition of maps corresponds to a product of Jacobians,
in the same order as the composition --- and since matrix multiplication is not
commutative, the order is now load-bearing.

\begin{remark}
A network is exactly such a composition, alternating affine maps and elementwise
activations. Its Jacobian is therefore an alternating product
\(W_L D_{L-1} W_{L-1}\cdots D_1 W_1\) with \(D_k\) diagonal. The vanishing and
exploding behaviour anticipated in \refsec{long-compositions} is now a statement
about the singular values of that product.
\end{remark}

\section{The Hessian}
\labsec{hessian}

\begin{definition}
For \(f:\R^{n}\to\R\) with continuous second partials, the \emph{Hessian} is
\(H\in\R^{n\times n}\) with \(H_{ij}=\partial^{2}f/\partial x_i\partial x_j\).
\end{definition}

Equality of mixed partials makes \(H\) symmetric, so the spectral theorem of
\refsec{spectral} applies: it has real eigenvalues and an orthonormal
eigenbasis. Those eigenvalues are the curvatures of \(f\) along the
corresponding directions.

\section{Taylor approximation}
\labsec{taylor}

Extending \refeq{first-order} one order further,
\begin{equation}
	f(\vect{x}+\vect{h})=f(\vect{x})+\nabla f(\vect{x})\tp\vect{h}
	+\tfrac12\,\vect{h}\tp H(\vect{x})\vect{h}+o(\lVert\vect{h}\rVert^{2}).
	\labeq{taylor2}
\end{equation}

The quadratic term is exactly the quadratic form of \refsec{quadratic-forms}, so
everything proved there transfers. In particular a positive definite Hessian
makes the local model a bowl with a unique minimum.

\begin{proposition}[Second-derivative test]
Let \(\nabla f(\vect{x})=\vect{0}\). If \(H(\vect{x})\) is positive definite the
point is a local minimum; if negative definite, a local maximum; if it has
eigenvalues of both signs, a saddle point.
\end{proposition}

\begin{example}
For \(f(x_1,x_2)=x_1^{2}+3x_1x_2+2x_2^{2}\),
\(H=\begin{bmatrix}2&3\\3&4\end{bmatrix}\) with \(\det H=-1<0\), so the
eigenvalues have opposite signs and the origin is a saddle.
\end{example}

\marginnote{In high dimensions saddles vastly outnumber local minima, since a
critical point is a minimum only if \emph{every} one of \(n\) eigenvalues is
positive. Much of the practical behaviour of gradient descent on deep models is
about escaping saddles rather than escaping minima.}

\section{The Jacobian determinant}
\labsec{jacobian-determinant}

When \(m=n\), \(|\det J|\) measures how the map scales volume locally --- the
statement anticipated at the end of \refch{determinants}. It appears whenever
variables are changed inside an integral:
\[
	\int_{\vect{g}(\Omega)}f(\vect{y})\,d\vect{y}
	=\int_{\Omega}f\bigl(\vect{g}(\vect{x})\bigr)\,\bigl|\det J_{\vect{g}}(\vect{x})\bigr|\,d\vect{x}.
\]

\begin{remark}
This is the one place Part~II needs integration, and it is needed for
\refch{random-variables}: transforming a random variable transforms its density
by exactly this factor. A map with \(\det J=0\) collapses volume, which is the
calculus reading of singularity.
\end{remark}
